\documentclass[11pt]{article}
\usepackage[a4paper,margin=1in]{geometry}
\usepackage[T1]{fontenc}
\usepackage{lmodern}
\usepackage{amsmath,amssymb,amsthm,booktabs,enumitem,mathtools,microtype}
\usepackage[numbers,sort&compress]{natbib}
\usepackage[hidelinks]{hyperref}

\newtheorem{theorem}{Theorem}[section]
\newtheorem{proposition}[theorem]{Proposition}
\newtheorem{lemma}[theorem]{Lemma}
\newtheorem{corollary}[theorem]{Corollary}
\theoremstyle{definition}
\newtheorem{definition}[theorem]{Definition}
\theoremstyle{remark}
\newtheorem{remark}[theorem]{Remark}

\newcommand{\eps}{\varepsilon}
\newcommand{\PhiP}{\Phi^{P}}
\newcommand{\PhiJ}{\Phi^{J}}
\newcommand{\PhiI}{\Phi^{I}}
\newcommand{\PsiK}{\Psi^{[K]}}
\newcommand{\GapP}{\operatorname{Gap}_{P}}
\newcommand{\GapJ}{\operatorname{Gap}_{J}}
\newcommand{\GapI}{\operatorname{Gap}_{I}}
\newcommand{\GapK}{\operatorname{Gap}_{K}}
\newcommand{\Binfty}{B_{\infty}}

\hypersetup{
  pdftitle={Rank-One P2-Scale Prime-Tail Resummation and Bounded-Gap Necessity},
  pdfauthor={RH research program},
  pdfsubject={Uniform factorial resummation below the P2 scale with a bounded-prime-gap necessity theorem},
  pdfkeywords={prime gaps, prime-square tails, asymptotic series, Euler products, P2 scale}
}

\title{Rank-One $P_2$-Scale Prime-Tail Resummation\\
and Bounded-Gap Necessity}
\author{RH research program}
\date{August 8, 2026}

\begin{document}
\maketitle

\begin{abstract}
Let $x=p_y$, $L=\log x\ge512$, and retain the exact first strict prime
tail $P_1(y)$.  For every integer $1\le K\le\lfloor3L\rfloor$, we
replace all higher tails by the factorial Laplace kernel
\[
 I_{2r}^{[K]}=\frac{x^{1-2r}}{(2r-1)L}
 \sum_{j=0}^{K-1}\frac{(-1)^j j!}{((2r-1)L)^j}\qquad(r\ge2).
\]
For the seven endpoint channels and the exact RH-383 endpoint map, the
resulting rank-one surrogate $\GapK$ satisfies
\[
 \pi^2|\GapP-\GapK|
 \le \frac{x^{-3}}{L}\left(
 7560\eps_x+\frac{1638}{x^2}
 +1176\frac{K!}{(3L)^K}\right),
\]
where $\eps_x=0.027L^{1.801}\exp[-0.1853L^{3/5}(\log L)^{-1/5}]$.
The factorial ratio is at most $1/(3L)$ throughout the full integer
window, so the maximum over all such $K$ is $o(P_2(y))$.

This rank-one retention is sharp within the frozen $P/J/I$ kernel
hierarchy.  Maynard's theorem on infinitely many bounded consecutive
prime gaps implies
\[
 \limsup_y p_y^2|P_1(y)-I_2(p_y)|\ge\frac12
\]
and, with $X_\infty>0$ from the exact endpoint,
\[
 \limsup_y p_y^2\pi^2|\GapP-\GapI|\ge X_\infty,
 \qquad
 \limsup_y p_y^2\pi^2|\GapP-\GapJ|\ge X_\infty.
\]
Thus a fully smooth $J/I$ surrogate cannot attain the $P_2$ scale,
whereas preserving the exact $P_1$ coordinate does.  The result is not a
claim that the factorial series converges and not a universal
impossibility theorem for every conceivable surrogate.  A 56-row exact
artifact reproduces the algebraic interfaces but is not an analytic
proof.
\end{abstract}

\noindent\textbf{Keywords:} bounded consecutive prime gaps; strict prime
tails; factorial asymptotics; rank-one retention; square-clock endpoint;
$P_2$ scale.

\section{Setting, inherited endpoint, and main results}

Let $3=p_1<p_2<\cdots$ be the odd primes and put
\[
 x=p_y,\qquad L=\log x,\qquad
 V(L)=L^{3/5}(\log L)^{-1/5},\qquad
 \eps_x=\frac{27}{1000}L^{1801/1000}
 e^{-(1853/10000)V(L)}.
\]
Johnston and Yang's Theorem~1.4, equation~(1.8), printed page~2, gives
\begin{equation}
 |\vartheta(t)-t|\le t\eps_t\qquad(t\ge23),
 \label{eq:JY}
\end{equation}
and RH-386 proves that the displayed envelope decreases once
$\log t\ge512$ \cite{JohnstonYang2023,RH386}.  Hence $\eps_t\le\eps_x$
for $t\ge x$ in our domain.

For $r\ge1$ define the strict prime and integral tails
\begin{equation}
 P_r(y)=\sum_{p>x}(p^2-1)^{-r},\qquad
 J_r(x)=\int_x^\infty\frac{dt}{(t^2-1)^r\log t},\qquad
 I_{2r}(x)=\int_x^\infty\frac{t^{-2r}}{\log t}\,dt.
 \label{eq:tails}
\end{equation}
For $c=1,\ldots,7$ let
\begin{equation}
 \PhiP_c=\sum_{r\ge1}\frac{c^rP_r}{r},\qquad
 \PhiJ_c=\sum_{r\ge1}\frac{c^rJ_r}{r},\qquad
 \PhiI_c=\sum_{r\ge1}\frac{c^rI_{2r}}{r}.
 \label{eq:PhiPJI}
\end{equation}
The nonnegative Tonelli identities identify $\PhiJ_c$ and $\PhiI_c$
with the logarithmic integrals in RH-387 \cite{RH387}.

For $r\ge2$, set
\begin{equation}
 K_r=\frac{x^{1-2r}}{(2r-1)L},\qquad
 a_r=\frac{1}{(2r-1)L},\qquad
 S_K(a)=\sum_{j=0}^{K-1}(-1)^j j!a^j,
 \label{eq:factorial-defs}
\end{equation}
and
\begin{equation}
 I_{2r}^{[K]}=K_rS_K(a_r),\qquad
 \PsiK_c=cP_1(y)+\sum_{r\ge2}\frac{c^rI_{2r}^{[K]}}r.
 \label{eq:Psi}
\end{equation}
The first coordinate is therefore retained exactly; every rank $r\ge2$
is replaced.

We recall the exact endpoint map.  For $2\le m\le8$, put
\[
 u_m=\prod_{p\ \mathrm{odd}}
 \frac{1-m/p^2}{1-1/p^2},
\]
and index
\[
 (\alpha_m)=(-2,2,-2,2,-2,2,-2),\qquad
 (\beta_m)=(1,-2,2,-2,2,-2,2).
\]
With
\[
 C(Z)=1+\sum_{m=2}^8\alpha_mZ_m,
 \qquad W(Z)=\sum_{m=2}^8\beta_mZ_m,
 \qquad Z_m(z)=u_me^{z_{m-1}},
\]
define
\begin{equation}
 F(z)=2\{C(u)-C(Z(z))\}
 -4W(Z(z))(1-e^{-z_1}).
 \label{eq:F}
\end{equation}
The exact normal form of RH-383 \cite{RH383} gives, for
$q_y=4\prod_{i\le y}p_i^2$,
\begin{equation}
 \GapP:=\Binfty-G(q_y)=\frac{F(\PhiP)}{\pi^2},\quad
 \GapJ:=\frac{F(\PhiJ)}{\pi^2},\quad
 \GapI:=\frac{F(\PhiI)}{\pi^2},\quad
 \GapK:=\frac{F(\PsiK)}{\pi^2}.
 \label{eq:gaps}
\end{equation}
Here each undecorated coordinate symbol denotes its seven-vector.  Only
$\GapP$ is the actual square-clock gap; the other three are declared
surrogates.

\begin{theorem}[Uniform rank-one $P_2$-scale resummation]
\label{thm:sufficiency}
Let $x=p_y$, $L\ge512$, and let $K$ be any integer satisfying
$1\le K\le\lfloor3L\rfloor$.  The vectors $\PhiP,\PhiJ,\PhiI,$ and
$\PsiK$ belong to $[0,1/2]^7$, and
\begin{equation}
 \max_{1\le c\le7}|\PhiP_c-\PsiK_c|
 \le\frac{x^{-3}}L\left(
 60\eps_x+\frac{13}{x^2}
 +\frac{28}{3}\frac{K!}{(3L)^K}\right).
 \label{eq:coordinate-master}
\end{equation}
Consequently,
\begin{equation}
 \boxed{
 \pi^2|\GapP-\GapK|
 \le\frac{x^{-3}}L\left(
 7560\eps_x+\frac{1638}{x^2}
 +1176\frac{K!}{(3L)^K}\right).}
 \label{eq:gap-master}
\end{equation}
Moreover,
\begin{equation}
 \lim_{y\to\infty}
 \max_{1\le K\le\lfloor3\log p_y\rfloor}
 \frac{|\GapP-\GapK|}{P_2(y)}=0.
 \label{eq:uniform-P2}
\end{equation}
\end{theorem}

Define the positive endpoint constant
\begin{equation}
 X_\infty=2u_4-4u_5+6u_6-8u_7+10u_8>0.
 \label{eq:X}
\end{equation}
Its positivity is the exact run-count consequence proved in RH-381
\cite{RH381}.

\begin{theorem}[Bounded-gap necessity in the frozen hierarchy]
\label{thm:necessity}
Within the $P/J/I$ hierarchy defined above,
\begin{align}
 \limsup_{y\to\infty}p_y^2|P_1(y)-I_2(p_y)|&\ge\frac12,
 \label{eq:scalar-I}\\
 \limsup_{y\to\infty}p_y^2|P_1(y)-J_1(p_y)|&\ge\frac12,
 \label{eq:scalar-J}\\
 \limsup_{y\to\infty}p_y^2\pi^2|\GapP-\GapI|&\ge X_\infty,
 \label{eq:gap-I-necessity}\\
 \limsup_{y\to\infty}p_y^2\pi^2|\GapP-\GapJ|&\ge X_\infty.
 \label{eq:gap-J-necessity}
\end{align}
In particular,
\begin{equation}
 \limsup_y\frac{|\GapP-\GapI|}{P_2(y)}
 =\limsup_y\frac{|\GapP-\GapJ|}{P_2(y)}=+\infty.
 \label{eq:ratio-divergence}
\end{equation}
Thus the exact $P_1$ term is necessary for the specific smooth-surrogate
$P_2$ contract in Theorem~\ref{thm:sufficiency}.  No statement is made
about every conceivable surrogate.
\end{theorem}

\section{Higher-rank source and power transfers}

Write $E(t)=\vartheta(t)-t$ and
$h_r(t)=(t^2-1)^{-r}/\log t$.  The strict endpoint $p>x$ gives
\begin{equation}
 P_r-J_r=-E(x)h_r(x)-\int_x^\infty E(t)h_r'(t)\,dt,
 \qquad
 |P_r-J_r|\le\eps_x\{2xh_r(x)+J_r\}.
 \label{eq:strict-source}
\end{equation}
The boundary $-E(x)h_r(x)$ and the second copy of $xh_r(x)$ produced by
integration by parts are distinct.  We apply \eqref{eq:strict-source}
only for $r\ge2$.

Let
\begin{equation}
 R(z)=-\log(1-z)-z=\sum_{r\ge2}\frac{z^r}{r}.
 \label{eq:R}
\end{equation}
For $0\le z<1$,
\begin{equation}
 0\le R(z)\le\frac{z^2}{2(1-z)},
 \qquad R'(z)=\frac{z}{1-z}.
 \label{eq:R-bounds}
\end{equation}
The first estimate retains the divisor $2$; the second will control the
change from $(t^2-1)^{-1}$ to $t^{-2}$.

\begin{proposition}[Higher-rank coordinate ledger]
\label{prop:coordinates}
For every $c\in\{1,\ldots,7\}$ and $L\ge512$,
\begin{align}
 \left|\sum_{r\ge2}\frac{c^r(P_r-J_r)}r\right|
 &<\frac{60\eps_x}{x^3L},
 \label{eq:source60}\\
 0\le\sum_{r\ge2}\frac{c^r(J_r-I_{2r})}r
 &<\frac{13}{x^5L}.
 \label{eq:power13}
\end{align}
\end{proposition}

\begin{proof}
Since $L\ge512$,
\begin{equation}
 x=e^L>2^L\ge2^{512}>256.
 \label{eq:bridge}
\end{equation}
All summands are absolutely summable, so Tonelli may be applied to their
nonnegative majorants.  Summing the boundary term in
\eqref{eq:strict-source} and using \eqref{eq:R-bounds} gives
\[
 2x\sum_{r\ge2}\frac{c^rh_r(x)}r
 \le\frac{x^{-3}}L
 \frac{c^2}{(1-x^{-2})^2\{1-c/(x^2-1)\}}.
\]
The $J_r$ contribution obeys
\[
 \sum_{r\ge2}\frac{c^rJ_r}r
 =\int_x^\infty\frac{R(c/(t^2-1))}{\log t}\,dt
 \le\frac{x^{-3}}L
 \frac{c^2}{6(1-x^{-2})^2\{1-c/(x^2-1)\}}.
\]
The reciprocal denominator common to these two terms and to the power
comparison below is
\begin{equation}
 D_c(x)^{-1}:=
 \frac1{(1-x^{-2})^2\{1-c/(x^2-1)\}}
 =\frac1{(1-x^{-2})\{1-(1+c)/x^2\}}.
 \label{eq:Dc}
\end{equation}
It increases with $c$ and decreases with $x$, so
\begin{equation}
 D_c(x)^{-1}<D_7(256)^{-1}
 =\frac{536870912}{536797185}<\frac{36}{35}<\frac54.
 \label{eq:Dc-bound}
\end{equation}
The source coefficient is therefore
\[
 \frac76c^2D_c(x)^{-1}<\frac65c^2
 \le\frac{294}{5}<60,
\]
which proves \eqref{eq:source60} by explicit arithmetic.

For the second comparison, the mean-value theorem and
$R'(z)=z/(1-z)$ give
\[
 0\le R\!\left(\frac{c}{t^2-1}\right)-R\!\left(\frac{c}{t^2}\right)
 \le\frac{c^2}{t^2(t^2-1)(t^2-1-c)}.
\]
Therefore
\[
 \sum_{r\ge2}\frac{c^r(J_r-I_{2r})}r
 \le\frac{x^{-5}}L
 \frac{c^2}{5(1-x^{-2})\{1-(1+c)/x^2\}}.
\]
By \eqref{eq:Dc-bound}, the last coefficient is
\[
 \frac{c^2}{5}D_c(x)^{-1}<\frac{c^2}{4}
 \le\frac{49}{4}<13.
\]
This proves \eqref{eq:power13} by explicit arithmetic.
\end{proof}

\section{Exact factorial remainder, cube, and moving window}

The power kernel has a Laplace representation with an exact finite
remainder.

\begin{lemma}[Factorial Laplace identity]
\label{lem:factorial}
For $a>0$, let
\[
 G(a)=\int_0^\infty\frac{e^{-v}}{1+av}\,dv.
\]
Then $I_{2r}=K_rG(a_r)$ and, for every integer $K\ge1$,
\begin{equation}
 G(a)-S_K(a)=(-a)^K
 \int_0^\infty\frac{e^{-v}v^K}{1+av}\,dv.
 \label{eq:exact-remainder}
\end{equation}
Consequently
\begin{equation}
 (-1)^K(I_{2r}-I_{2r}^{[K]})\ge0,\qquad
 |I_{2r}-I_{2r}^{[K]}|\le K_r a_r^K K!.
 \label{eq:factorial-bound}
\end{equation}
\end{lemma}

\begin{proof}
Substitute $t=xe^u$ in $I_{2r}$ and then set
$v=(2r-1)u$.  This gives $I_{2r}=K_rG(a_r)$.  The finite geometric
identity
\[
 \frac1{1+av}=\sum_{j=0}^{K-1}(-av)^j
 +\frac{(-av)^K}{1+av}
\]
may be integrated termwise because it is finite.  Since
$\int_0^\infty e^{-v}v^j\,dv=j!$, it gives
\eqref{eq:exact-remainder}; positivity of its integral and
$1/(1+av)\le1$ give \eqref{eq:factorial-bound}.
\end{proof}

\begin{proposition}[Factorial coordinate and cube]
\label{prop:factorial-coordinate}
If $1\le K\le\lfloor3L\rfloor$, then
\begin{equation}
 \max_{1\le c\le7}
 \left|\sum_{r\ge2}\frac{c^r(I_{2r}-I_{2r}^{[K]})}r\right|
 \le\frac{28}{3}\frac{x^{-3}}L\frac{K!}{(3L)^K}.
 \label{eq:factorial28}
\end{equation}
Moreover $0<S_K(a_r)\le1$ for every $r\ge2$, and $\PsiK$ belongs to
$[0,1/2]^7$.
\end{proposition}

\begin{proof}
Lemma~\ref{lem:factorial}, $2r-1\ge3$, and $1/r\le1/2$ give
\begin{align*}
 \sum_{r\ge2}\frac{c^r}{r}|I_{2r}-I_{2r}^{[K]}|
 &\le \frac{x^{-3}}L\frac{K!}{(3L)^K}
 \frac{c^2}{6(1-c/x^2)}\\
 &<\frac{4c}{3}\frac{x^{-3}}L\frac{K!}{(3L)^K}.
\end{align*}
Maximizing $4c/3$ over $c\le7$ proves \eqref{eq:factorial28}.  This
sum is over all $r\ge2$, not only $r=2$.

The magnitudes $j!a_r^j$ in $S_K(a_r)$ have successive ratio
$(j+1)a_r$.  Within the stated window every ratio that occurs is at most
$(K-1)/(3L)<1$.  Pairing the decreasing alternating terms, with the
case $K=1$ immediate, gives $0<S_K(a_r)\le1$.

Finally,
\[
 P_1(y)\le\sum_{n>x}\frac1{n^2-1}
 =\frac12\left(\frac1x+\frac1{x+1}\right),
\]
whereas
\[
 0\le\sum_{r\ge2}\frac{c^rI_{2r}^{[K]}}r
 \le\frac{c^2}{6x^3L(1-c/x^2)}.
\]
At $x>256$, $L\ge512$, their sum after multiplication of the first term
by $c\le7$ is below $1/2$.  This proves the cube assertion.
\end{proof}

\begin{lemma}[Uniform integer window]
\label{lem:window}
For every integer $1\le K\le\lfloor3L\rfloor$,
\begin{equation}
 b_K:=\frac{K!}{(3L)^K}\le b_1=\frac1{3L}.
 \label{eq:bK}
\end{equation}
\end{lemma}

\begin{proof}
For $K+1\le\lfloor3L\rfloor$,
\[
 \frac{b_{K+1}}{b_K}=\frac{K+1}{3L}\le1.
\]
Induction from $b_1$ proves the assertion over the complete integer
window.  Finite checks of small $K$ are only regression fixtures and are
not used as a proof of this quantified statement.
\end{proof}

\section{Endpoint derivative ledgers and proof of sufficiency}

The factors in $u_m$ have summable deficits, and the $p=3$ factor gives
\begin{equation}
 0<u_m\le\frac{9-m}{8}\qquad(2\le m\le8).
 \label{eq:u-bound}
\end{equation}
Thus the endpoint arrays satisfy
\begin{equation}
 A:=\sum_{m=2}^8|\alpha_m|u_m\le7,\qquad
 B:=\sum_{m=2}^8|\beta_m|u_m\le\frac{49}{8}.
 \label{eq:AB}
\end{equation}

\begin{lemma}[Endpoint gradient and Hessian]
\label{lem:endpoint-derivatives}
For every $z\in[0,1/2]^7$,
\begin{align}
 \|\nabla F(z)\|_1&<126,
 \label{eq:gradient126}\\
 \sum_{i,j=1}^7|\partial_{ij}F(z)|&<224.
 \label{eq:Hessian224}
\end{align}
In particular,
\begin{equation}
 |F(z)-\nabla F(0)\mathbin{\cdot}z|
 \le112\|z\|_\infty^2.
 \label{eq:Taylor112}
\end{equation}
\end{lemma}

\begin{proof}
For the gradient, the derivative of the $C$ term, the derivative falling
on $W$, and the derivative falling on $1-e^{-z_1}$ have coefficients
$2,4,4$.  Since $e^{z_j}\le e^{1/2}<2$, their entrywise ledger is
\[
 \|\nabla F(z)\|_1
 \le e^{1/2}(2A+4B+4B)
 \le63e^{1/2}<126.
\]
This is the dual $\ell^1$ gradient for an $\ell^\infty$ input.

For the Hessian, the $C$ term contributes $2A$ before the exponential
factor.  In the product of $W$ and the loss $1-e^{-z_1}$, the Hessian of
$W$, the two gradient cross terms, and the Hessian of the loss contribute
$4B,8B,4B$.  Hence
\[
 \sum_{i,j}|\partial_{ij}F(z)|
 \le e^{1/2}\{2A+(4+8+4)B\}
 \le112e^{1/2}<224.
\]
Taylor's theorem on the segment from $0$ to $z$ gives
\eqref{eq:Taylor112}, including the factor $1/2$.
\end{proof}

\begin{proof}[Proof of Theorem~\ref{thm:sufficiency}]
The cube assertion for $\PhiP,\PhiJ,\PhiI$ is the direct integer-tail and
integral estimate of RH-387; Proposition~\ref{prop:factorial-coordinate}
proves it for $\PsiK$.  Insert between $\PhiP_c$ and $\PsiK_c$ the two
coordinates
\[
 cP_1+\sum_{r\ge2}\frac{c^rJ_r}{r},
 \qquad
 cP_1+\sum_{r\ge2}\frac{c^rI_{2r}}{r}.
\]
Propositions~\ref{prop:coordinates} and
\ref{prop:factorial-coordinate} give \eqref{eq:coordinate-master}.

The segment from $\PhiP$ to $\PsiK$ lies in the convex cube.  The
$\ell^\infty/\ell^1$ mean-value estimate in
Lemma~\ref{lem:endpoint-derivatives} therefore multiplies the three
coordinate constants by $126$:
\[
 126\cdot60=7560,\qquad
 126\cdot13=1638,\qquad
 126\cdot\frac{28}{3}=1176.
\]
Together with \eqref{eq:gaps}, this proves \eqref{eq:gap-master}.

RH-384 proves the intrinsic scale
\begin{equation}
 P_2(y)\sim\frac1{3x^3L}
 \label{eq:P2-scale}
\end{equation}
\cite{RH384}.  Lemma~\ref{lem:window} makes the right side of
\eqref{eq:gap-master} uniform in all admissible $K$.  Since
$\eps_x\to0$, $x^{-2}\to0$, and $L^{-1}\to0$, division by
\eqref{eq:P2-scale} proves \eqref{eq:uniform-P2}.
\end{proof}

\section{Bounded consecutive gaps and the rank-one obstruction}

Maynard's unconditional Theorem~1.3, printed page~385 (PDF page~3),
states
\begin{equation}
 \liminf_{n\to\infty}(p_{n+1}-p_n)\le600
 \label{eq:Maynard}
\end{equation}
\cite{Maynard2015}.  We use only the resulting infinitely many
\emph{consecutive} prime pairs $x=p_y$, $q=p_{y+1}=x+h$ with
$h\le600$.  Indeed the gaps are integers: if only finitely many were at
most $600$, every sufficiently late gap would be at least $601$, forcing
the liminf to be at least $601$.

\begin{lemma}[Successor jump]
\label{lem:jump}
Let $E_y=P_1(y)-I_2(p_y)$.  Along the bounded consecutive gaps supplied
by \eqref{eq:Maynard},
\begin{equation}
 E_y-E_{y+1}=\frac1{q^2-1}
 -\int_x^q\frac{dt}{t^2\log t},
 \qquad x^2(E_y-E_{y+1})\longrightarrow1.
 \label{eq:E-jump}
\end{equation}
Consequently $\limsup_y p_y^2|E_y|\ge1/2$.
\end{lemma}

\begin{proof}
The exact strict successor identity is
$P_1(y)=(q^2-1)^{-1}+P_1(y+1)$.  Subtract the two integral tails to get
the first formula.  Since $h\le600$, $q/x\to1$ and
\[
 \frac{x^2}{q^2-1}\to1,
 \qquad
 0\le x^2\int_x^q\frac{dt}{t^2\log t}
 \le\frac{h}{\log x}\to0.
\]
Thus the scaled jump tends to one.  The triangle inequality at the two
endpoints, together with $q/x\to1$, shows that at least one endpoint of
each such pair has its own normalized absolute value at least
$1/2-o(1)$.  Infinitely many pairs supply infinitely many distinct
endpoints, proving the limsup assertion.  An eventual $1/16$ lower
witness is a safe finite certificate row, not the sharp constant.
\end{proof}

The same conclusion holds with $J_1$ because
\begin{equation}
 0\le J_1(x)-I_2(x)
 =\int_x^\infty\frac{dt}{t^2(t^2-1)\log t}
 =O\!\left(\frac1{x^3L}\right)=o(x^{-2}).
 \label{eq:J1-I2}
\end{equation}

We now lift the scalar jump through the full endpoint.  Set
$D_y=\PhiP(y)-\PhiI(p_y)$.  Exact successors and Tonelli give, in channel
$c$,
\begin{equation}
 (D_y-D_{y+1})_c
 =-\log\!\left(1-\frac{c}{q^2-1}\right)
 -\int_x^q\frac{-\log(1-c/t^2)}{\log t}\,dt.
 \label{eq:vector-jump}
\end{equation}
Uniformly for $c=1,\ldots,7$, the first term is
$c/q^2+O(q^{-4})$, while the scaled integral is $O(h/L)$.  Hence, along
the same bounded gaps,
\begin{equation}
 x^2(D_y-D_{y+1})\longrightarrow(1,2,3,4,5,6,7).
 \label{eq:rank-one-direction}
\end{equation}

The endpoint direction is exact.  From \eqref{eq:F},
\begin{align}
 \nabla F(0)\mathbin{\cdot}(1,2,3,4,5,6,7)
 &=\sum_{m=2}^8u_m\{-2\alpha_m(m-1)-4\beta_m\}\\
 &=4u_4-8u_5+12u_6-16u_7+20u_8
 =2X_\infty.
 \label{eq:gradient-direction}
\end{align}
The individual coefficients for $m=2,\ldots,8$ are
$0,0,4,-8,12,-16,20$.

\begin{proof}[Proof of Theorem~\ref{thm:necessity}]
Lemma~\ref{lem:jump} proves \eqref{eq:scalar-I}, and
\eqref{eq:J1-I2} proves \eqref{eq:scalar-J}.  The direct integral bounds
of RH-387 give $\PhiI=O(1/(xL))$.  Its source comparison gives
$\PhiP=O(1/(xL))$ as well.  Lemma~\ref{lem:endpoint-derivatives} therefore
implies
\begin{align*}
 F(\PhiP)-F(\PhiI)
 &=\nabla F(0)\mathbin{\cdot}(\PhiP-\PhiI)+o(x^{-2}),
\end{align*}
because $x^2\|\PhiP\|_\infty^2+x^2\|\PhiI\|_\infty^2=O(L^{-2})$.

Equations \eqref{eq:rank-one-direction} and
\eqref{eq:gradient-direction} show that the scalar sequence
$\nabla F(0)\cdot D_y$ has scaled successor jump tending to
$2X_\infty$.  The same two-endpoint triangle argument as in
Lemma~\ref{lem:jump} yields limsup at least $X_\infty$.  The Taylor
remainder is $o(x^{-2})$, so \eqref{eq:gaps} proves
\eqref{eq:gap-I-necessity}.

RH-387 gives
$\pi^2|\GapJ-\GapI|\le588/(x^3L)=o(x^{-2})$, which transfers the same
lower bound to \eqref{eq:gap-J-necessity}.  Finally divide these
subsequence lower bounds by \eqref{eq:P2-scale}; since $xL\to\infty$,
the two limsup ratios are infinite.  This proves the theorem.
\end{proof}

\section{Novelty, limitations, and claim firewall}

RH-387 replaced every rank, including $P_1$, by a smooth integral.  Its
source error $\eps_x/(xL)$ is larger than the $P_2$ scale because
$\eps_xx^2\to\infty$.  The present theorem does not sharpen that source
estimate.  Instead it keeps $P_1$ exact and applies the source transfer
only to $r\ge2$, gaining the factor $x^{-2}$ needed for a uniform
$P_2$-scale result.  The exact Laplace remainder simultaneously compiles
every $r\ge2$ and every integer $K$ in a moving window; this is not a
finite-partition corollary of RH-386 or RH-387.

The bounded-gap theorem supplies the complementary negative edge:
within the declared $P/J/I$ hierarchy, smoothing the rank-one coordinate
creates an actual-prime jump of order $x^{-2}$, much larger than $P_2$.
This establishes necessity for this hierarchy only.  It does not exclude
surrogates carrying different exact prime data or other correction terms.

The finite sums $S_K$ are asymptotic truncations.  We claim neither that
$\sum_j(-1)^jj!a^j$ converges nor that the window may be enlarged beyond
$K\le\lfloor3L\rfloor$.  The paid error gives no $P_3$ or cubic
coefficient precision.  Channels are the seven real integers; there is
no complex-$c$ theorem.  The finite clock is fixed before the prefix
limit, $c_{11}=0$ remains phasewise, and neither the clock nor an order
$K_N$ grows with the prefix length.  No active-$c_{11}$ cancellation,
operator, trace, prime-power trace, zeta-zero identification, or proof of
the Riemann hypothesis is supplied.  Gates A--E are all false.

\section{Executable reproduction and source integrity}
\label{sec:artifact}

The exact artifact has epistemic role
\texttt{reproduction\_not\_analytic\_proof}.  Its 56 rows comprise 12
analytic interfaces, seven channel ledgers, 12 finite factorial
regressions, seven endpoint rows, ten bounded-gap/Taylor rows, and eight
master ledgers.  The canonical certificate has 14,531 bytes and SHA-256
\begin{center}
\small\ttfamily
373d870847bb0bf134aa1eba30c5e4d2c3a01dba470af9c75ebacadd81976371.
\end{center}
It independently recomputes the exact fractions, the universal
$b_K$ induction, the $60,13,28/3$ coordinate constants, the
$126,224,112$ endpoint ledgers, and the $7560,1638,1176$ products.
Twenty-four genuine semantic mutations are rejected without calling the
fresh certificate or any row builder.  Strict JSON parsing, exact Python
types, optimized \texttt{-OO} replay, and a closed official Draft
2020--12 schema are tested.  These finite checks do not prove
\eqref{eq:JY}, Maynard's theorem, Stieltjes integration, Tonelli's
theorem, or the asymptotic limits.

The immutable Git closure contains 77 blobs from the RH-387 release
commit
\begin{center}
\small\ttfamily
dedd8e8d2c44564e66524a646f9cf5fb9a389c77.
\end{center}
Two canonical remote objects, Johnston--Yang and Maynard, give 79 logical
sources.  The ordered
Git and logical digests are
\begin{center}
\small\ttfamily
\begin{tabular}{@{}l@{}}
d7f2ee43f56631c8f3442db8fcc6fb423a801b5af7607351623cd449a92c3f73\\
bffce602d6e3b568eb96662820f08aa457ff5d0de4065f3c9eeac53d8d8dfa39.
\end{tabular}
\end{center}
The canonical remote-lock digests, in source-key order, are
\begin{center}
\small\ttfamily
\begin{tabular}{@{}l@{}}
d53b93212b7c5b5b6b3f7e890099c48ce8e35f2bff9bdd49f9c330a9b5039786\\
bd4aad4b7042218e5733bb07db2a513770710628a8ac52d5bcc9881fcb0b5d2e.
\end{tabular}
\end{center}

Neither external PDF nor the Johnston--Yang source tar is vendored or
archived.  Network verification is disabled by default and requires an
explicit opt-in.  Offline semantic verification freezes the Maynard
article page and DOI.  The opt-in retrieval gate checks only the exact
requested/final PDF URL, status, PDF MIME type, 528,115 bytes, 31 pages,
and SHA-256.  The version of record is Copyright 2015 Department of
Mathematics, Princeton University.  Current publisher policy materials
do not establish the agreement applicable to this 2015 article, so the
release conservatively records no article-specific redistribution grant
and \texttt{redistributable\_in\_release=false}.  The inherited
Johnston--Yang author-manuscript lock remains nonredistributable as well.

\section*{Declarations}

\paragraph{Data and code availability.}
The proof, certificate compiler, strict tests, closed schema, two remote
source-lock records, offline-by-default verification, and release replay
tools accompany this paper.  The external PDFs and source tar are not
redistributed.

\paragraph{Author contributions.}
The author is responsible for conceptualization, proof, software,
validation, writing, and release auditing.

\paragraph{Funding.}
No external funding was received.

\paragraph{Competing interests.}
The author declares no competing interests.

\paragraph{Ethics.}
No human participants, animals, personal data, or clinical interventions
are involved.

\paragraph{AI assistance.}
AI-assisted tools were used for exact symbolic-interface enumeration,
adversarial testing, manuscript drafting, and release auditing.  All
mathematical claims, citations, source locks, and final text remain under
author responsibility.

\bibliographystyle{plainnat}
\bibliography{references}
\end{document}
